%%%%%%%%%%  scan idx 314  |  folio 301  %%%%%%%%%%
\(H(K[X_1,\ldots,X_n])\) d'après (1.3.ii)).

Remarquons enfin que si \(H=\Ghat\), \(u_H^n(1+Xt)=\prod\limits_{i=1}^n(1+T_it)=1+X_1t+\cdots+X_nt^n\).

\textbf{Théorème 1.6.} L'application de \(\Hom_{\Abfil}(\Ghat_K,H)\) dans \(H(K[X])\) qui, à un morphisme de foncteurs, associe sa valeur sur l'élément \(1+Xt\) de \(\Ghat_K(K[X])\), est injective; son image est le sous-groupe de \(H(K[X])\) formé des éléments \(\theta\) tels que (avec les notations de 1.4. pour les \(f_n\))~:
\[
\begin{array}{ll}
\text{i)}& H(f_1)(\theta)=0.\\[0.3em]
\text{ii)}& \forall n\geq1,\quad
H(f_n)\bigl(u_H^n(\theta)\bigr)\equiv u_H^n(\theta)
\ \mathrm{mod.}\ F^nH(K[X_1,\ldots,X_n]).
\end{array}
\]

Montrons d'abord l'injectivité. D'après le lemme 1.2, il suffit de voir qu'un morphisme de foncteurs \(\Phi\) nul sur \(1+Xt\) est nul sur \(\varphi=1+X_1t+\cdots+X_nt^n+\cdots\). Le diagramme commutatif~:
\[
\begin{tikzcd}
\Ghat(K[X]) \arrow[r,"u^n_{\Ghat}"] \arrow[d,"\Phi"'] &
\Ghat(K[X_1,\ldots,X_n]) \arrow[d,"\Phi"]\\
H(K[X]) \arrow[r,"u_H^n"] &
H(K[X_1,\ldots,X_n])
\end{tikzcd}
\tag{1.6.1}
\]
montre que pour tout \(n\), \(\Phi(1+X_1t+\cdots+X_nt^n)=0\).

Or \(\varphi\equiv 1+X_1t+\cdots+X_nt^n\ \mathrm{mod.}\ F^{n+1}\Ghat\)~; donc \(\Phi(\varphi)\equiv \Phi(1+\cdots+X_nt^n)\ \mathrm{mod.}\ F^{n+1}H\), et par suite \(\Phi(\varphi)\equiv0\ \mathrm{mod.}\ F^{n+1}H\), pour tout \(n\)~; comme \(H(K[X_i]_{i\geq 1})\) est séparé et complet, \(\Phi(\varphi)=0\).

Soit maintenant \(\Phi\) un élément quelconque de \(\Hom(\Ghat,H)\), et soit 
